\subsection{\texorpdfstring{第~\ref{sec:penrose-singularity-theorem}~节}{Penrose 奇性定理}}\label{subsec:solutions-penrose-singularity-theorem}

\begin{description}[font=\normalfont,leftmargin=2em,labelwidth=1.5em]
\RGsolution{ex:christoffel-symbols-in-null-coordinates}{christoffel-symbols-in-null-coordinates}
容易看出
\[
  (g^{\mu\nu})=
  \begin{pmatrix}
    \alpha&-1&\beta_2&\beta_3\\
    -1&0&0&0\\
    \beta_2&0&\gamma_{22}&\gamma_{23}\\
    \beta_3&0&\gamma_{32}&\gamma_{33}
  \end{pmatrix}^{-1}
  =
  \begin{pmatrix}
    0&-1&0&0\\
    -1&\delta&\beta^2&\beta^3\\
    0&\beta^2&\gamma^{22}&\gamma^{23}\\
    0&\beta^3&\gamma^{32}&\gamma^{33}
  \end{pmatrix},
\]
其中 $\beta^i=\sum_{j=2}^3\gamma^{ij}\beta_j$, $\delta=\sum_{i=2}^3\beta_i\beta^i-\alpha$. 因此, 例如有
\begin{align*}
  \Gamma_{vr}^v
   &=\frac{1}{2}\sum_{\alpha=0}^3g^{v\alpha}
      \left(\frac{\partial g_{r\alpha}}{\partial v}
            +\frac{\partial g_{v\alpha}}{\partial r}
            -\frac{\partial g_{vr}}{\partial x^\alpha}\right)\\
   &=-\frac{1}{2}\left(\frac{\partial g_{rr}}{\partial v}
            +\frac{\partial g_{vr}}{\partial r}
            -\frac{\partial g_{vr}}{\partial r}\right)=0,
\end{align*}
同理, $\Gamma_{rr}^v=\Gamma_{ri}^v=\Gamma_{rr}^r=\Gamma_{rr}^i=0$. 最后,
\begin{align*}
  \Gamma_{rj}^i
   &=\frac{1}{2}\sum_{\alpha=0}^3g^{i\alpha}
       \left(\frac{\partial g_{j\alpha}}{\partial r}
             +\frac{\partial g_{r\alpha}}{\partial x^j}
             -\frac{\partial g_{rj}}{\partial x^\alpha}\right)\\
   &=\frac{1}{2}\beta^i
       \left(\frac{\partial g_{vj}}{\partial r}
             +\frac{\partial g_{vr}}{\partial x^j}
             -\frac{\partial g_{rj}}{\partial v}\right)\\
   &\quad+\frac{1}{2}\sum_{k=2}^3\gamma^{ik}
       \left(\frac{\partial g_{jk}}{\partial r}
             +\frac{\partial g_{rk}}{\partial x^j}
             -\frac{\partial g_{rj}}{\partial x^k}\right)\\
   &=\frac{1}{2}\sum_{k=2}^3\gamma^{ik}\frac{\partial\gamma_{jk}}{\partial r}
       =\sum_{k=2}^3\gamma^{ik}\beta_{kj}.
\end{align*}

\RGsolution{ex:limits-of-penrose-singularity-theorem}{limits-of-penrose-singularity-theorem}
\begin{enumerate}[label=(\alph*),ref=\alph*,leftmargin=*]
\item Minkowski 时空不包含陷俘曲面.
\item Einstein 宇宙具有紧 Cauchy 超曲面.
\item de Sitter 宇宙具有紧 Cauchy 超曲面.
\item 二维反 de Sitter 时空不是整体双曲的; 但应注意, 该时空的任何整体双曲开子集都不是测地完备的.
\end{enumerate}
\end{description}
